\subsection{From the dual numerator to bags of balls}\label{hb:section}

We give the geometric construction with explicit losses.  No boundary
ratio or connectivity assumption is used in this construction; those
assumptions enter the subsequent charging estimates.  Throughout this
subsection, $\mathcal T$ is a hierarchy of a finite graph $G=(V,E)$,
$X_v\in\{0,1\}^d$, and $U:\mathbb R^d\to\mathbb R^E$ is a linear
contraction.  For an oriented edge $e=uv$, put
\[
 X_e=X_u-X_v,\qquad a_e=\langle u_e,X_e\rangle,\qquad
 d_e=\|X_e\|_2^2=\|X_e\|_1.
\]
Here $u_e^{\mathsf T}$ is row $e$ of $U$.  We assume $a_e\ge0$; multiplying individual
rows by signs preserves contraction and replaces each $a_e$ by $|a_e|$,
which can only increase every squared boundary sum under consideration.
In particular $a_e^2\le d_e$.  An edge is counted with its multiplicity
and occurs in at most two sets $O(t)$.

\begin{definition}[Balls, bags, and geometric sequences]\label{hb:definitions}
Set $Q=[0,1]^d$ and
\[
 B_Q(x,r)=\{y\in Q:\|y-x\|_1<r\}\qquad(r>0).
\]
A \emph{bag} is a nonempty collection of pairwise disjoint such balls,
all of the same radius and with centers among the $X_v$.  Its type is
$(\delta)$ when that radius is $\delta$, and its type is $(\delta,\Delta)$
when, in addition, all distances between its centers are at most
$\Delta$.  Its size $|\mathrm{Bag}|$ is its number of balls.
For $\beta>1$, a bag of type $(\delta,\Delta)$ is
\emph{$\beta$-compact} if
\[
 |\mathrm{Bag}|\ge2,\qquad
 \beta\Delta\le\delta|\mathrm{Bag}|.
\]
The midpoint of two centers belongs to $Q$, so disjoint radius-$\delta$
balls have centers at distance at least $2\delta$.  Hence a compact bag
has $\Delta\ge2\delta$.

A bag of radius $\delta$ is \emph{$\eta$-assigned to $t$} if
\[
 |\mathrm{Bag}_t|\ge\eta|O(t)|,
\]
and every ball has a center $X_u$ with $u\in V(t)$ and an edge
$uv\in O(t)$ such that $d_{uv}<\delta$.  A \emph{family} of bags is
nonempty, and all balls in all its bags are mutually disjoint.
A compact family has type $(\delta,\Delta)$ when its bags have that
common type.  An assigned family has type $(\delta,T)$ when its bags
are assigned to distinct nodes in $T$ and have radius $\delta$.

For $0<\lambda<1$, a sequence of compact families of types
$(\delta_i,\Delta_i)$ is \emph{$\lambda$-geometric} if
$\Delta_{i+1}<\lambda\delta_i$ for all successive families.
A sequence of assigned families of types $(\delta_i,T_i)$ is
$\lambda$-geometric if $\delta_{i+1}<\lambda\delta_i$ and the sets
$T_i$ are pairwise disjoint.  Its \emph{mass} is
\[
 \mathsf M=\sum_i\delta_i
       \sum_{\mathrm{Bag}\in\mathcal F_i}|\mathrm{Bag}|.
\]
\end{definition}

\subsubsection{Spectral packing}
\label{hb:spectral-subsection}

We first prove the geometric estimate underlying the construction of
bags.  It is the spectral packing statement of
\cite[Lemma~6.4]{AOG}; the argument below gives a logarithmic bound and
chooses the radius explicitly.
Here a Euclidean ball of \emph{squared radius} \(r\) means
\[
 \mathcal B_2(y,r)=\{z:\|z-y\|_2^2\le r\}.
\]
In particular, two such closed balls are disjoint when their centers
are at distance strictly greater than \(2\sqrt r\).

\begin{lemma}[Spectral packing]\label{hb:spectral-packing}
Let \(F\) be a nonempty set of oriented edges of a finite multigraph
with edge set \(E\), and put \(m=|F|\).
For points \(Y_v\in\R^E\), write
\(Y_e=Y_u-Y_v\) when \(e=(u,v)\), and denote the \(e\)-coordinate of
\(Y_e\) by \(Y_{e,e}\).
Suppose that \(a>0\) and
\begin{equation}\label{hb:spectral-assumption}
 \Upsilon:=\left(\frac1m\sum_{e\in F}Y_{e,e}\right)^2>0,
 \qquad
 \Upsilon\ge\frac a m\sum_{e\in F}\|Y_e\|_2^2.
\end{equation}
Then \(a\le1\), and there are \(b\ge2\) pairwise disjoint Euclidean
balls of one positive squared radius \(r\), centered at images of
endpoints of edges in \(F\), such that
\begin{equation}\label{hb:spectral-logarithmic}
 b\ge\frac{am}{16},
 \qquad
 br\ge\frac{9m\Upsilon}
 {512\bigl(5+\log_2(1/a)\bigr)^2}.
\end{equation}
Consequently, for every \(0<\varepsilon<1/3\), the choice
\begin{equation}\label{hb:spectral-constant}
 C_1(\varepsilon)=\frac{1024}{\varepsilon^2}
\end{equation}
gives
\begin{equation}\label{hb:spectral-conclusion}
 b\ge\frac{a|F|}{C_1(\varepsilon)},
 \qquad
 br\ge\frac{a^\varepsilon\Upsilon|F|}{C_1(\varepsilon)}.
\end{equation}
\end{lemma}

\begin{proof}
The scalar Cauchy--Schwarz inequality gives
\[
 \Upsilon\le\frac1m\sum_{e\in F}Y_{e,e}^2
 \le\frac1m\sum_{e\in F}\|Y_e\|_2^2.
\]
The last quantity is positive, so the assumption implies \(a\le1\).

\smallskip
\noindent\emph{The matrix estimate.}
Let \(P_F:\R^E\to\R^F\) be coordinate restriction, and form the
\(m\)-by-\(m\) matrix \(A\) whose column indexed by \(e\in F\)
is \(P_FY_e\).  Set \(\tau=m\sqrt\Upsilon\).
The singular values \(\sigma_1\ge\cdots\ge\sigma_m\ge0\) of \(A\)
satisfy
\begin{equation}\label{hb:spectral-matrix-bounds}
 \sum_{i=1}^m\sigma_i\ge|\tr A|=\tau,
 \qquad
 \sum_{i=1}^m\sigma_i^2=\|A\|_F^2
 \le\frac{\tau^2}{am}.
\end{equation}
Indeed, writing a singular value decomposition as
\(A=\sum_i\sigma_i u_i v_i^{\mathsf T}\), the trace is
\(\sum_i\sigma_i\langle u_i,v_i\rangle\), whose absolute value
is at most \(\sum_i\sigma_i\).  The Frobenius norm identity follows
from orthonormality, and its upper bound follows from
\eqref{hb:spectral-assumption} and the fact that coordinate restriction
does not increase a vector's norm.

We will also use the following rank approximation inequality:
for \(0\le k\le m\) and any matrix \(C\) of rank at most \(k\),
\begin{equation}\label{hb:spectral-rank-bound}
 \|A-C\|_F^2\ge\sum_{i=k+1}^m\sigma_i^2.
\end{equation}
For completeness, let \(Q\) be orthogonal projection onto the column
space of \(C\).  Since \((I-Q)C=0\), orthogonal projection and the
singular value decomposition give
\[
 \|A-C\|_F^2
 \ge\|(I-Q)A\|_F^2
 =\sum_{i=1}^m\sigma_i^2\bigl(1-\|Qu_i\|_2^2\bigr).
\]
The numbers \(w_i=\|Qu_i\|_2^2\) lie in \([0,1]\), and
\(\sum_iw_i=\rank Q\le k\).
Because the \(\sigma_i^2\) are nonincreasing,
\(\sum_i\sigma_i^2w_i\le\sum_{i=1}^k\sigma_i^2\):
moving available weight from a later index to an earlier index never
decreases this weighted sum, and its maximum therefore puts weight one
on the first \(k\) indices.  This proves
\eqref{hb:spectral-rank-bound}, including \(k=0\).

\smallskip
\noindent\emph{Selecting a singular-value tail.}
Define
\[
 p_0=\left\lfloor\frac{am}{16}\right\rfloor+1,
 \qquad T_p=\sum_{i=p}^m\sigma_i^2\quad(1\le p\le m).
\]
We have \(1\le p_0\le m\).  By \eqref{hb:spectral-matrix-bounds},
\[
 \sum_{i<p_0}\sigma_i
 \le\sqrt{(p_0-1)\sum_i\sigma_i^2}
 \le\frac\tau4.
\]
The sum is zero when \(p_0=1\), so this also covers arbitrarily small
\(am\).  It follows that
\begin{equation}\label{hb:spectral-surviving-sum}
 \sum_{i=p_0}^m\sigma_i\ge\frac{3\tau}{4}.
\end{equation}
Partition the integer interval \([p_0,m]\) into the nonempty blocks
\[
 I_j=\{2^jp_0,\ldots,\min(2^{j+1}p_0-1,m)\},
 \qquad 0\le j<L,
\]
where
\[
 L=1+\left\lfloor\log_2\frac m{p_0}\right\rfloor
 \le 5+\log_2(1/a).
\]
For the first index \(p=2^jp_0\) of a block, its cardinality is at
most \(p\).  Cauchy--Schwarz therefore gives
\[
 \sum_{i\in I_j}\sigma_i\le\sqrt{pT_p}.
\]
Summing this inequality and using
\eqref{hb:spectral-surviving-sum}, one of these first indices \(p\)
satisfies
\begin{equation}\label{hb:spectral-large-tail}
 p\ge p_0>\frac{am}{16},
 \qquad
 pT_p\ge\frac{9\tau^2}{16L^2}>0.
\end{equation}

\smallskip
\noindent\emph{Selecting the balls at an explicit radius.}
Fix the positive number
\begin{equation}\label{hb:spectral-radius}
 r=\frac{T_p}{32m}.
\end{equation}
Scan the endpoint images of edges in \(F\), keeping a point precisely
when its distance from every previously kept point is greater than
\(2\sqrt r\).  This finite procedure produces centers
\(Y_{w_1},\ldots,Y_{w_b}\) for pairwise disjoint closed balls of
squared radius \(r\).  Its maximality also ensures that each endpoint
image \(Y_v\) lies within distance \(2\sqrt r\) of one of the
centers; choose such a center and call its vertex \(c(v)\).

If \(b\le p\), define a matrix \(C\) by
\[
 C_e=P_F\bigl(Y_{c(u)}-Y_{c(v)}\bigr)
 \qquad(e=(u,v)\in F).
\]
Every column of \(C\) belongs to the span of
\(P_F(Y_{w_j}-Y_{w_1})\), \(2\le j\le b\), so
\(\rank C\le b-1\le p-1\).  On the other hand,
\begin{align*}
 \|A-C\|_F^2
 &\le\sum_{e=(u,v)\in F}
 \bigl\|(Y_u-Y_{c(u)})-(Y_v-Y_{c(v)})\bigr\|_2^2\\
 &\le\sum_{e=(u,v)\in F}
 2\bigl(\|Y_u-Y_{c(u)}\|_2^2+
          \|Y_v-Y_{c(v)}\|_2^2\bigr)
 \le16mr=\frac{T_p}{2}.
\end{align*}
This contradicts \eqref{hb:spectral-rank-bound} with \(k=p-1\),
which bounds the same quantity below by \(T_p>0\).
Thus \(b\ge p+1\ge2\).  Combining
\eqref{hb:spectral-large-tail} and \eqref{hb:spectral-radius} yields
\[
 b>\frac{am}{16},
 \qquad
 br\ge\frac{pT_p}{32m}
 \ge\frac{9m\Upsilon}{512L^2}.
\]
The bound on \(L\) proves \eqref{hb:spectral-logarithmic}.

\smallskip
\noindent\emph{Conversion to the power bound.}
Put \(t=\log(1/a)\ge0\).  Since \(\log2\ge1/2\),
\(5+\log_2(1/a)\le5+2t\).
The elementary maximum of \(t e^{-\varepsilon t/2}\) is
\(2/(e\varepsilon)\le1/\varepsilon\), and
\(5\le5/(3\varepsilon)\) for \(0<\varepsilon<1/3\).
Consequently,
\[
 \bigl(5+\log_2(1/a)\bigr)e^{-\varepsilon t/2}
 \le5+\frac2\varepsilon<\frac4\varepsilon.
\]
Squaring and substituting in \eqref{hb:spectral-logarithmic} gives
\[
 br\ge\frac{9\varepsilon^2}{8192}
       a^\varepsilon m\Upsilon
 \ge\frac{\varepsilon^2}{1024}
       a^\varepsilon m\Upsilon.
\]
Finally \(C_1(\varepsilon)\ge16\), so the count estimate in
\eqref{hb:spectral-conclusion} follows as well.
\end{proof}

The positivity assumption on \(\Upsilon\) is needed when one asks for
many distinct centers.  If \(\Upsilon=0\),
\eqref{hb:spectral-assumption} forces \(Y_e=0\) for every \(e\in F\),
and all endpoint images may coincide.  Such a set contributes zero to
the squared sums for which this lemma is used and can be discarded.


\subsubsection{Homogeneity and the two constructions}

\begin{definition}[Bad nodes and homogeneous subsets]\label{hb:bad}
For $0<\alpha\le1$, a nonroot node $t$ with $O(t)\ne\varnothing$
is \emph{$\alpha$-bad} if
\[
 \left(\frac{\sum_{e\in O(t)}a_e}{|O(t)|}\right)^2
 \ge\frac{\alpha}{|O(t)|}\sum_{e\in O(t)}d_e.
\]
Only nodes with a positive boundary sum need be retained.
A set of edges is \emph{$c$-homogeneous} if its $a_e^2$ are positive
and any two of them have ratio less than $c$, and the same holds for
its $d_e$.  A subset $O'(t)\subseteq O(t)$ is
\emph{$\gamma$-dominating} if
\[
 \left(\sum_{e\in O'(t)}a_e\right)^2
 \ge\gamma\left(\sum_{e\in O(t)}a_e\right)^2.
\]
\end{definition}

\begin{lemma}[Two elementary truncations]\label{hb:bucketing}
Define
\[
 J=1+\left\lceil\log_2(256/\alpha^2)\right\rceil,
 \qquad \gamma=\frac1{4J^4}.
\]
Every $\alpha$-bad node with a positive boundary sum admits a
$2$-homogeneous, $\gamma$-dominating subset $O'(t)$.
\end{lemma}
\begin{proof}
Fix the node, let $m=|O(t)|$, and let expectations be uniform on $O(t)$.
Writing $\mu=\mathbb E a_e>0$, badness says
$\mathbb E d_e\le\mu^2/\alpha$.  Discard edges for which
$a_e<\mu/4$ or $d_e>16\mu^2/\alpha^2$.  The first class contributes
at most $\mu/4$ to $\mathbb E a_e$.  For the second class, Markov's
inequality gives probability at most $\alpha/16$, and therefore
Cauchy--Schwarz and $a_e^2\le d_e$ give
\[
 \mathbb E\bigl[a_e\mathbf1_{\{d_e>16\mu^2/\alpha^2\}}\bigr]
 \le\sqrt{\mathbb E a_e^2\,
          \mathbb P\{d_e>16\mu^2/\alpha^2\}}
 \le\mu/4.
\]
The remaining edges contribute at least $m\mu/2$ to the boundary sum.
On those edges both $a_e^2$ and $d_e$ lie in
$[\mu^2/16,16\mu^2/\alpha^2]$.  Divide this interval into half-open
dyadic intervals, with an extra last interval if needed.  There are
at most $J$ choices for each quantity and at most $J^2$ pairs of
choices.  Some pair contains edges whose $a_e$ sum is at least
$m\mu/(2J^2)$.  Its edges are $2$-homogeneous, and squaring this
last estimate gives the asserted dominance.
\end{proof}

\begin{lemma}[One homogeneous family]\label{hb:homogeneous-bags}
Let $T$ be a nonempty set of $\alpha$-bad nodes with positive boundary
sums, and let $0<\gamma\le1$.  Suppose that $O'(t)$ is
$\gamma$-dominating for each $t\in T$
and that $F=\bigcup_{t\in T}O'(t)$ is $4$-homogeneous.  Write
\[
 s=\frac{\alpha\gamma}{32},\qquad
 c_1=\min_{e\in F}a_e^2,
 \qquad
 S'=\sum_{t\in T}\frac1{|O(t)|}
                     \left(\sum_{e\in O'(t)}a_e\right)^2.
\]
For $\beta>1$ and $0<\epsilon<1/3$, put
$C=C_1(\epsilon)=1024\epsilon^{-2}$.  If
\begin{equation}\label{hb:smallness}
 s^\epsilon\le\frac1{48\beta C},
\end{equation}
there is either a family of $\beta$-compact bags or a family of
$s^{1+2\epsilon}$-assigned bags, assigned to a subset of $T$, whose
mass is at least
\begin{equation}\label{hb:one-mass}
 \frac{s^\epsilon}{192\beta C}\,S'.
\end{equation}
Its radius, and also its diameter parameter in the compact case,
belong to $[sc_1,12c_1/s]$.
\end{lemma}
\begin{proof}
Set $m_t=|O(t)|$, $m'_t=|O'(t)|$,
$N=\sum_{t\in T}m_t$, $N'=|F|$, and $c_2=\max_{e\in F}d_e$.
The sets $O'(t)$ are nonempty.  Homogeneity, dominance, and badness
imply, for each $t$,
\[
 4c_1(m'_t)^2
 \ge\left(\sum_{O'(t)}a_e\right)^2
 \ge\gamma\left(\sum_{O(t)}a_e\right)^2
 \ge\alpha\gamma m_t\sum_{O(t)}d_e
 \ge\frac{\alpha\gamma}{4}\,m_tm'_tc_2.
\]
Cancel $m'_t$ and sum.  Since
$\sum_t m'_t\le2N'$, we obtain
\begin{equation}\label{hb:density}
 c_1N'\ge sNc_2.
\end{equation}
Also $N\ge N'$ and $c_1\le c_2$, since $a_e^2\le d_e$.
Consequently
\[
 s\le\widetilde\alpha:=sN/N'\le c_1/c_2\le1.
\]
Let $Y_v=UX_v$.  The row-$e$ coordinate of $Y_u-Y_v$ is $a_e$.
Moreover $\|Y_u-Y_v\|_2^2\le d_e$, so
\[
 (\mathbb E_{e\in F}a_e)^2
 \ge c_1\ge\widetilde\alpha c_2
 \ge\widetilde\alpha\mathbb E_{e\in F}\|Y_e\|_2^2.
\]
Lemma~\ref{hb:spectral-packing} yields $b$ disjoint Euclidean balls
with squared radius $r$ and centers at endpoint images $Y_v$, with
\[
 b\ge\widetilde\alpha N'/C,
 \qquad br\ge\widetilde\alpha^\epsilon c_1N'/C.
\]
Their center distances squared are at least $4r$.  Contraction and
the binary-coordinate identity therefore give $\|X_u-X_v\|_1\ge4r$
for distinct selected centers.  Thus the corresponding $Q$-balls of
radius $r$ are disjoint.  Shrink their common radius to $\delta$ so that
\begin{equation}\label{hb:normalized-mass}
 b\delta=w:=\widetilde\alpha^\epsilon c_1N'/C.
\end{equation}
Call this family of $b$ small balls $\mathcal A$.
There are at most $2N'$ distinct endpoint centers.  It follows that
\[
 \frac{s^\epsilon c_1}{2C}\le\delta
 \le c_1\widetilde\alpha^{\epsilon-1}\le c_1/s.
\]
As $\epsilon<1/2$, condition~\eqref{hb:smallness} implies
$s^{1-\epsilon}\le s^\epsilon\le1/(2C)$; hence
\begin{equation}\label{hb:radius-range}
 sc_1\le\delta\le c_1/s.
\end{equation}
Finally, homogeneity and the multiplicity bound show
\begin{equation}\label{hb:numerator-mass}
 S'\le4c_1\sum_t\frac{(m'_t)^2}{m_t}
 \le8c_1N'\le8Cs^{-\epsilon}w.
\end{equation}

Let $V'(t)$ consist of the endpoints of $O'(t)$ that lie in $V(t)$,
and let $V'=\bigcup_tV'(t)$.  Define
$D=\max\{\delta,2c_2\}$.  Choose a maximal disjoint collection
$\mathcal B$ of radius-$D$ balls centered at $X_u$ for $u\in V'$.
Every such center is within $2D$ of a chosen center: otherwise its ball
could be added.  Each center of a small ball is an endpoint of an edge
of $F$, so it is within $c_2$ of a point of $V'$, and hence within $3D$
of a chosen center.

If $|\mathcal B|<w/(12\beta D)$, assign every small ball to its closest
chosen center, breaking ties consistently.  The balls assigned to one
center have center diameter at most $6D$.  Keep a resulting bag
precisely when it contains at least $6\beta D/\delta$ balls.
The discarded bags contain fewer than
$|\mathcal B|6\beta D/\delta<b/2$ balls.  Each retained bag is
$\beta$-compact of type $(\delta,6D)$; the numerical threshold is
larger than $2$ because $D\ge\delta$ and $\beta>1$.
Their balls are still mutually disjoint and their mass is at least
$w/2$, which implies~\eqref{hb:one-mass}.

In the other case, $|\mathcal B|\ge w/(12\beta D)$, assign each ball
of $\mathcal B$ to one node $t$ for which its center belongs to $V'(t)$.
Keep the bag assigned to $t$ when
\[
 |\mathrm{Bag}_t|\ge\frac{|\mathcal B|m_t}{2N}.
\]
Discarding the other bags loses fewer than $|\mathcal B|/2$ balls.
For each retained ball its center is in $V(t)$ and some edge of
$O'(t)$ has length at most $c_2<D$.  Furthermore
\begin{equation}\label{hb:assignment-size}
 \frac{|\mathrm{Bag}_t|}{m_t}
 \ge\frac{w}{24\beta DN}.
\end{equation}
If $D=\delta$, the right side is $b/(24\beta N)$, at least
$s/(24\beta C)$ by the bound on $b$.  If $D=2c_2$, it is at least
$s^{1+\epsilon}/(48\beta C)$ by
\eqref{hb:density} and~\eqref{hb:normalized-mass}.
In either case~\eqref{hb:smallness} makes it at least
$s^{1+2\epsilon}$.  These bags are therefore assigned as asserted.
Their mass is at least $D|\mathcal B|/2\ge w/(24\beta)$;
\eqref{hb:numerator-mass} gives~\eqref{hb:one-mass}.

In both cases~\eqref{hb:density} gives $c_2\le c_1/s$ and therefore
$D\le2c_1/s$.  Together with~\eqref{hb:radius-range}, this places all
radii, and the compact diameter $6D$, in $[sc_1,12c_1/s]$.
\end{proof}

\Needspace{18\baselineskip}
\subsubsection{Selecting separated scales}

\begin{proposition}[Geometric bag extraction]\label{prop:hierarchy-bags}
Let $T$ be a nonempty finite set of $\alpha$-bad nodes with positive boundary
sums, where $0<\alpha\le1$.  Define
\[
 \begin{aligned}
 J&=1+\lceil\log_2(256/\alpha^2)\rceil,
 &C_2(\alpha)&=128J^4,
 &s&=\alpha/C_2(\alpha),\\
 L&=2+\lceil\log_2(1/s)\rceil,
 &H&=1+\lceil\log_2(24/(\lambda s^2))\rceil,
 &C_1(\epsilon)&=1024\epsilon^{-2}.
 \end{aligned}
\]
Suppose $\beta>1$, $0<\epsilon<1/3$, $0<\lambda<1$, and
$s^\epsilon\le(48\beta C_1(\epsilon))^{-1}$.
There is a $\lambda$-geometric sequence of families, consisting either
entirely of $\beta$-compact bags or entirely of
$s^{1+2\epsilon}$-assigned bags assigned to distinct nodes of $T$,
whose mass satisfies
\begin{equation}\label{hb:extraction-bound}
 \mathsf M\ge
 \frac{s^\epsilon}{384\beta C_1(\epsilon)C_2(\alpha)LH}
 \sum_{t\in T}\frac1{|O(t)|}
                  \left(\sum_{e\in O(t)}a_e\right)^2.
\end{equation}
In particular, all losses are functions of $\alpha,\epsilon,\beta$
and $\lambda$, independent of the number of vertices and hierarchy depth.
\end{proposition}
\begin{proof}
Lemma~\ref{hb:bucketing} gives $2$-homogeneous subsets $O'(t)$ with
$\gamma=1/(4J^4)=32/C_2(\alpha)$.  Define
\[
 A_t=\min_{e\in O'(t)}a_e^2,
 \qquad B_t=\min_{e\in O'(t)}d_e,
 \qquad s_t=\frac{(\sum_{O'(t)}a_e)^2}{|O(t)|}.
\]
All are positive.  For $m=|O(t)|$ and $m'=|O'(t)|$, the same
pointwise argument used above, now with $2$-homogeneity, gives
\[
 2A_t(m')^2\ge\gamma\left(\sum_{O(t)}a_e\right)^2
 \ge\alpha\gamma mm'B_t.
\]
Consequently
\begin{equation}\label{hb:local-scale-ratio}
 1\le B_t/A_t\le2/(\alpha\gamma)=1/(16s).
\end{equation}
Partition the nodes according to the two integers
$n=\lfloor\log_2 A_t\rfloor$ and
$m=\lfloor\log_2 B_t\rfloor$.  For each fixed $n$,
\eqref{hb:local-scale-ratio} allows at most $L$ values of $m$.
Retain the value $m(n)$ maximizing the sum of the $s_t$ over that group,
and call the resulting node set $T_n$.  Summing over $n$ yields
\begin{equation}\label{hb:retained-sum}
 \sum_n\sum_{t\in T_n}s_t
 \ge\frac1L\sum_{t\in T}s_t
 \ge\frac1{C_2(\alpha)L}
       \sum_{t\in T}\frac{(\sum_{O(t)}a_e)^2}{|O(t)|}.
\end{equation}
The last inequality uses $\gamma\ge1/C_2(\alpha)$.

The union of the $O'(t)$ over $T_n$ is $4$-homogeneous: its $a_e^2$
lie in $[2^n,2^{n+2})$ and its $d_e$ in
$[2^{m(n)},2^{m(n)+2})$.  Apply
Lemma~\ref{hb:homogeneous-bags} to every nonempty $T_n$.
Let $\mathcal F_n$ be the resulting family.  Its radius is at least
$s2^n$ and at most $24\,2^n/s$; the latter upper bound also holds for
its compact diameter parameter.  Its mass is at least
$s^\epsilon\sum_{t\in T_n}s_t/(192\beta C_1(\epsilon))$.

Partition these families by their integer $n$ modulo $H$ and by whether
they are compact or assigned.  One of the $2H$ resulting classes has
at least $1/(2H)$ of their total mass.  Order that class by decreasing
$n$.  If $n'$ follows $n$, then $n'\le n-H$, and the definition of $H$
implies
\[
 \frac{24\,2^{n'}}s
 \le\frac{24\,2^{n-H}}s
 <\lambda s2^n.
\]
This verifies the geometric condition, including the stronger diameter
condition for compact families.  Assigned nodes are distinct across the
sequence because the original $T_n$ are disjoint and each construction
uses only a subset of its own $T_n$.  Finally combine its $1/(2H)$ mass
share with~\eqref{hb:retained-sum} and the mass guarantee for each
$\mathcal F_n$.  This is exactly~\eqref{hb:extraction-bound}.
\end{proof}
